% chapters/bab5-kesimpulan-dan-saran.tex -- BAB 5 KESIMPULAN DAN SARAN
% Contoh teks ilmiah di bawah ini adalah placeholder untuk memperlihatkan
% hasil format (thesis.md); ganti dengan isi penelitian yang sebenarnya.

\chapter{KESIMPULAN DAN SARAN}

\section{Kesimpulan}
Berdasarkan hasil penelitian dan pembahasan pada bab sebelumnya, dapat
disimpulkan beberapa hal sebagai berikut.
\begin{enumerate}
  \item Tahapan praproses yang disesuaikan dengan karakteristik
        bahasa Indonesia, khususnya normalisasi kata tidak baku,
        terbukti memberikan kontribusi yang signifikan terhadap
        kualitas data sebelum digunakan pada model klasifikasi
        sentimen.
  \item Model LSTM dengan arsitektur satu lapisan tersembunyi berhasil
        mengklasifikasikan sentimen ulasan produk berbahasa Indonesia
        dengan kinerja yang lebih baik dibandingkan model
        \textit{baseline} berbasis SVM, dengan peningkatan F1-score
        sebesar 10 poin persentase pada data uji.
  \item Kemampuan LSTM dalam mempertimbangkan urutan kata terbukti
        efektif dalam menangani kalimat yang mengandung pola negasi,
        yang menjadi salah satu sumber kesalahan klasifikasi utama
        pada model berbasis fitur kata tunggal.
  \item Model yang diusulkan masih menunjukkan keterbatasan dalam
        menangani ulasan yang mengandung sarkasme atau opini implisit,
        sehingga tingkat kesalahan klasifikasi pada kategori tersebut
        relatif lebih tinggi dibandingkan kategori lainnya.
\end{enumerate}

\section{Saran}
Berdasarkan kesimpulan di atas, beberapa saran yang dapat
dipertimbangkan untuk penelitian selanjutnya adalah sebagai berikut.
\begin{enumerate}
  \item Penelitian selanjutnya dapat mempertimbangkan penggunaan
        model berbasis Transformer, seperti BERT yang telah
        dilatih khusus untuk bahasa Indonesia, guna meningkatkan
        kemampuan model dalam memahami konteks kalimat secara lebih
        mendalam.
  \item Penambahan jumlah dan variasi data latih, khususnya pada
        ulasan yang mengandung sarkasme, diharapkan dapat membantu
        model mempelajari pola opini implisit dengan lebih baik.
  \item Penelitian selanjutnya dapat mengeksplorasi kombinasi metode
        LSTM dengan mekanisme \textit{attention} untuk meningkatkan
        kemampuan model dalam berfokus pada bagian kalimat yang lebih
        relevan terhadap sentimen yang diekspresikan.
  \item Evaluasi lebih lanjut terhadap efisiensi komputasi model perlu
        dilakukan apabila model akan diterapkan pada sistem pemantauan
        ulasan produk secara \textit{real-time}.
\end{enumerate}
